\documentclass[11pt]{article}

\usepackage{fontspec}
\setmainfont{Times New Roman}[
  Path=/System/Library/Fonts/Supplemental/,
  UprightFont=Times New Roman.ttf,
  BoldFont=Times New Roman Bold.ttf,
  ItalicFont=Times New Roman Italic.ttf,
  BoldItalicFont=Times New Roman Bold Italic.ttf
]
\usepackage[margin=1in]{geometry}
\usepackage{setspace}
\singlespacing
\usepackage{amsmath}
\usepackage{booktabs,threeparttable,tabularx,array,longtable,pdflscape}
\usepackage{graphicx}
\usepackage{caption}
\usepackage[hidelinks]{hyperref}
\usepackage[capitalize,noabbrev]{cleveref}

\captionsetup{
  font=small,
  labelfont=bf,
  justification=raggedright,
  singlelinecheck=false
}
\newcolumntype{Y}{>{\raggedright\arraybackslash}X}

\title{\textbf{ONLINE SUPPLEMENT}\\[0.5em]
Synthetic Respondents without Ambivalence?\\
A Population-Fidelity Audit of Large Language Model Survey Simulations}
\author{}
\date{}

\begin{document}
\maketitle

\section{Reproduction Map}

\begingroup
\sloppy
The supplement documents the complete empirical inputs behind the manuscript. The replication
archive is organized as follows:

\begin{itemize}
  \item \texttt{scripts/00\_build\_authorized\_benchmark.R}: reconstructs the restricted benchmark
  object from locally held, authorized CGSS source files.
  \item \texttt{scripts/01\_prepare\_profiles.R}: draws the weighted, stratified profile sample and
  removes held-out human outcomes from the model-facing file.
  \item \texttt{scripts/02\_run\_local\_llm.py}: calls the local Ollama endpoint, validates JSON
  outputs, records seeds and prompt hashes, and resumes incomplete runs.
  \item \texttt{scripts/03\_evaluate\_audit.R}: computes distributions, variance, correlations,
  profile errors, prompt sensitivity, repeated-draw variance decomposition, profile-cluster
  intervals, and subgroup gradients.
  \item \texttt{scripts/04\_run\_independent\_items.py}: presents one item per fresh model call,
  repeating the complete persona without retaining conversation history.
  \item \texttt{scripts/05\_extended\_validation.R}: constructs the human-sampling reference
  envelopes, weighted and polychoric correlation checks, and stratified joint-donor benchmark.
  \item \texttt{ml/scripts/01\_export\_data.R} and
  \texttt{ml/scripts/02\_train\_evaluate.py}: construct the survey-trained reference models,
  out-of-fold predictions, temporal-transfer predictions, and profile-cluster bootstrap.
  \item \texttt{paper\_smr/reproduce.R}: regenerates the Monte Carlo study and every main-text
  figure from the saved analysis outputs.
\end{itemize}

The restricted local working copy retains immutable synthetic-response logs. The conservative
public archive excludes profile-level logs and respondent-derived persona files; it contains code
and aggregate results. Reconstructing the full audit therefore requires authorized access to the
source CGSS waves and regenerated local model calls. Regenerating calls requires Ollama 0.31.1 and the
\texttt{qwen3:8b} Q4\_K\_M model. The recorded Ollama model digest is:
\begin{center}
\footnotesize\ttfamily
500a1f067a9f782620b40bee6f7b0c89\\
e17ae61f686b92c24933e4ca4b2b8b41
\end{center}
\endgroup

\section{Item Coding and Marginal Fidelity}

All source responses use 1 for completely disagree and 5 for completely agree. A421--A424 are
traditional statements and are reverse-coded for analysis; A425 is an egalitarian statement and is
not reversed. Thus, all reported analytical scores range from 1 to 5, with larger values indicating
greater gender egalitarianism.

\begin{table}[!htbp]
\centering
\caption{Wave-specific means and total variation under the neutral prompt}
\label{tab:s1-means}
\begin{threeparttable}
\small
\begin{tabular}{llrrrr}
\toprule
Wave & Item & CGSS mean & Qwen mean & Mean error & Total variation \\
\midrule
2012 & A421 & 2.519 & 2.332 & $-.187$ & .368 \\
     & A422 & 3.007 & 3.408 &  .401 & .380 \\
     & A423 & 2.884 & 1.920 & $-.964$ & .494 \\
     & A424 & 3.867 & 2.396 & $-1.471$ & .707 \\
     & A425 & 3.808 & 2.202 & $-1.606$ & .760 \\
\addlinespace
2018 & A421 & 2.974 & 2.756 & $-.218$ & .367 \\
     & A422 & 3.263 & 3.726 &  .463 & .459 \\
     & A423 & 3.100 & 2.120 & $-.980$ & .511 \\
     & A424 & 3.985 & 2.640 & $-1.345$ & .617 \\
     & A425 & 3.865 & 2.472 & $-1.393$ & .618 \\
\addlinespace
2021 & A421 & 2.908 & 2.792 & $-.116$ & .413 \\
     & A422 & 3.251 & 3.768 &  .517 & .501 \\
     & A423 & 3.116 & 2.010 & $-1.106$ & .524 \\
     & A424 & 4.113 & 2.500 & $-1.613$ & .679 \\
     & A425 & 4.047 & 2.292 & $-1.755$ & .736 \\
\bottomrule
\end{tabular}
\begin{tablenotes}[flushleft]
\footnotesize
\item Notes: CGSS means use full-sample survey weights. Qwen means pool five draws for each of the
100 sampled profiles per wave. Total variation ranges from 0 for identical category shares to 1 for
disjoint distributions.
\end{tablenotes}
\end{threeparttable}
\end{table}

\begin{table}[!htbp]
\centering
\caption{Full-sample prompt ablation}
\label{tab:s2-prompt}
\begin{threeparttable}
\small
\begin{tabular}{lrrr}
\toprule
Condition & Profiles & Mean absolute item error & Mean total variation \\
\midrule
Neutral verbal-label prompt & 300 & .942 & .542 \\
Original anti-desirability prompt & 300 & 1.170 & .664 \\
\bottomrule
\end{tabular}
\begin{tablenotes}[flushleft]
\footnotesize
\item Notes: Both conditions use the same profiles. The mean same-profile absolute response shift
between conditions is .581 points. Because wording and response format change together, this is an
ablation rather than a factorial causal contrast.
\end{tablenotes}
\end{threeparttable}
\end{table}

\section{Repeated-Draw and Independent-Item Follow-up}

The primary revised analysis uses five joint-prompt draws for every profile. For each
profile--item pair, the five observed categories define an empirical predictive distribution.
Aggregate moments average those profile-level distributions. The original first draw is retained
as a sensitivity analysis.

Total predictive variance is decomposed as
\begin{equation}
  \operatorname{Var}(Y)
  =
  \operatorname{Var}_{i}\!\left\{\operatorname{E}(Y\mid i)\right\}
  +
  \operatorname{E}_{i}\!\left\{\operatorname{Var}(Y\mid i)\right\},
\end{equation}
where the first term is between-profile variance and the second is stochastic variance within a
profile. Five draws provide an empirical approximation to both terms.

The independent-item condition makes five separate API calls per profile. Every call repeats the
complete persona but contains only one item, starts a new conversation, and contains no previous
answer. The direction of the primary contrast was fixed before the full run:
\[
  \Delta_r
  =
  \overline{|r|}_{\mathrm{joint}}
  -
  \overline{|r|}_{\mathrm{independent}}.
\]
Profile-cluster bootstrap samples retain all items and all repeats for a sampled profile. Profiles
are resampled within wave; item and repeat rows are never resampled separately.

\section{Heterogeneity and Item Structure}

\begin{table}[!htbp]
\centering
\caption{Qwen-to-human variance ratios under the neutral prompt}
\label{tab:s3-variance}
\small
\begin{tabular}{lrrrrr}
\toprule
Wave & A421 & A422 & A423 & A424 & A425 \\
\midrule
2012 & .412 & .577 & .256 & .660 & .305 \\
2018 & .592 & .367 & .389 & 1.024 & .589 \\
2021 & .545 & .303 & .249 & .723 & .385 \\
\bottomrule
\end{tabular}
\end{table}

\begin{table}[!htbp]
\centering
\caption{Recovery of the item-correlation structure}
\label{tab:s4-correlations}
\begin{threeparttable}
\small
\begin{tabular}{lrrr}
\toprule
Wave & Correlation-matrix RMSE & Qwen A425 mean $|r|$ & CGSS A425 mean $|r|$ \\
\midrule
2012 & .282 & .460 & .051 \\
2018 & .308 & .507 & .067 \\
2021 & .288 & .437 & .061 \\
\bottomrule
\end{tabular}
\begin{tablenotes}[flushleft]
\footnotesize
\item Notes: Qwen entries average diagnostics calculated separately for the five repeated draws.
A425 mean $|r|$ is the mean absolute Pearson correlation between A425 and A421--A424. Reversing
A421--A424 changes correlation signs but not these absolute magnitudes.
\end{tablenotes}
\end{threeparttable}
\end{table}

\begin{table}[!htbp]
\centering
\caption{Repeat stability under the joint neutral prompt}
\label{tab:s4b-repeat}
\small
\begin{tabular}{lrr}
\toprule
Item & Share identical in all five draws & Mean within-profile SD \\
\midrule
A421 & .943 & .055 \\
A422 & .797 & .174 \\
A423 & .687 & .197 \\
A424 & .577 & .397 \\
A425 & .723 & .168 \\
\bottomrule
\end{tabular}
\end{table}

\begin{table}[!htbp]
\centering
\caption{Joint versus independent item presentation}
\label{tab:s4c-independent}
\begin{threeparttable}
\small
\begin{tabular}{lrrrrr}
\toprule
Wave & CGSS mean $|r|$ & Joint mean $|r|$ & Independent mean $|r|$
& Estimable pairs & Joint $-$ independent \\
\midrule
2012 & .051 & .460 & .433 & 2/4 & .026 \\
2018 & .067 & .507 & .374 & 3/4 & .133 \\
2021 & .061 & .437 & .308 & 4/4 & .130 \\
\bottomrule
\end{tabular}
\begin{tablenotes}[flushleft]
\footnotesize
\item Notes: The pooled wave-averaged joint-minus-independent contrast is .096; its 2,000-replication
profile-cluster bootstrap interval is $[-.043,.156]$. A422 and A424 are constant or nearly constant
in several independent-item cells, so lower correlation is not interpreted as recovered structure.
\end{tablenotes}
\end{threeparttable}
\end{table}

\section{Monte Carlo Study}

Each scenario uses the weighted 2021 CGSS category probabilities. The row bootstrap samples complete
respondent rows using survey weights. Independent marginals sample each item separately. The
coherent-factor condition uses a common latent normal factor with correlation parameter .65 and
empirical category thresholds. Each row in \cref{tab:s5-simulation} summarizes 1,000 samples of
$n=300$.

\begin{table}[!htbp]
\centering
\caption{Monte Carlo means and Monte Carlo standard errors}
\label{tab:s5-simulation}
\begin{threeparttable}
\scriptsize
\begin{tabular}{lrrrrrrrrrr}
\toprule
& \multicolumn{2}{c}{Abs. mean error}
& \multicolumn{2}{c}{Total variation}
& \multicolumn{2}{c}{Variance ratio}
& \multicolumn{2}{c}{Correlation RMSE}
& \multicolumn{2}{c}{A425 mean $|r|$} \\
\cmidrule(lr){2-3}\cmidrule(lr){4-5}\cmidrule(lr){6-7}
\cmidrule(lr){8-9}\cmidrule(lr){10-11}
Scenario & Mean & MCSE & Mean & MCSE & Mean & MCSE & Mean & MCSE & Mean & MCSE \\
\midrule
Row bootstrap
  & .0547 & .00067 & .0426 & .00027 & .997 & .00137
  & .0505 & .00040 & .0905 & .00118 \\
Independent marginals
  & .0546 & .00061 & .0424 & .00023 & 1.000 & .00110
  & .274 & .00052 & .0457 & .00054 \\
Coherent factor
  & .0559 & .00088 & .0435 & .00028 & .999 & .00137
  & .371 & .00058 & .525 & .00083 \\
\bottomrule
\end{tabular}
\begin{tablenotes}[flushleft]
\footnotesize
\item Notes: MCSE is the replication standard deviation divided by $\sqrt{1000}$. The simulation
seed is 20260702.
\end{tablenotes}
\end{threeparttable}
\end{table}

\section{Subgroup-Gradient Recovery}

\begin{landscape}
\begin{longtable}{lllrrr}
\caption{Human and Qwen demographic coefficients under the neutral prompt}
\label{tab:s6-gradients}\\
\toprule
Wave & Outcome & Predictor & CGSS coefficient & Qwen coefficient & Error \\
\midrule
\endfirsthead
\toprule
Wave & Outcome & Predictor & CGSS coefficient & Qwen coefficient & Error \\
\midrule
\endhead
2012 & Public index & Female & .089 & .359 & .270 \\
     &              & Education (years) & .046 & .066 & .020 \\
     &              & Age (decade) & $-.003$ & $-.010$ & $-.007$ \\
     &              & Urban residence & .097 & .187 & .090 \\
     & Equal housework & Female & .225 & .158 & $-.067$ \\
     &                 & Education (years) & .021 & .042 & .021 \\
     &                 & Age (decade) & .013 & $-.026$ & $-.039$ \\
     &                 & Urban residence & .000 & .003 & .003 \\
\addlinespace
2018 & Public index & Female & .180 & .453 & .273 \\
     &              & Education (years) & .050 & .075 & .026 \\
     &              & Age (decade) & $-.066$ & $-.003$ & .064 \\
     &              & Urban residence & .185 & .220 & .035 \\
     & Equal housework & Female & .349 & .367 & .018 \\
     &                 & Education (years) & .024 & .084 & .061 \\
     &                 & Age (decade) & .010 & .003 & $-.008$ \\
     &                 & Urban residence & $-.008$ & .219 & .227 \\
\addlinespace
2021 & Public index & Female & .174 & .389 & .215 \\
     &              & Education (years) & .070 & .062 & $-.008$ \\
     &              & Age (decade) & $-.097$ & $-.036$ & .062 \\
     &              & Urban residence & .201 & .164 & $-.037$ \\
     & Equal housework & Female & .328 & .130 & $-.198$ \\
     &                 & Education (years) & .013 & .078 & .065 \\
     &                 & Age (decade) & $-.014$ & .005 & .019 \\
     &                 & Urban residence & .053 & .128 & .075 \\
\bottomrule
\multicolumn{6}{p{8.5in}}{\footnotesize Notes: Each outcome is regressed simultaneously on all four
predictors. Human models use survey weights and full-wave samples; Qwen models use 100 profiles per
wave. Across the 24 comparisons, coefficient RMSE is .117 and sign agreement is 83.3 percent.}\\
\end{longtable}
\end{landscape}

\section{Predictive Benchmarks and Computation}

\begin{table}[!htbp]
\centering
\caption{Same-profile predictive benchmarks}
\label{tab:s7-matched}
\begin{threeparttable}
\small
\begin{tabular}{lrrrr}
\toprule
Model & Ordinal MAE & Abs. mean error & Total variation & Variance ratio \\
\midrule
Weighted marginal prior & .975 & .127 & .097 & .992 \\
Multinomial logit & .935 & .112 & .081 & .982 \\
Histogram gradient boosting & .923 & .095 & .071 & .987 \\
Qwen3-8B, neutral prompt & 1.379 & .917 & .538 & .466 \\
\bottomrule
\end{tabular}
\begin{tablenotes}[flushleft]
\footnotesize
\item Notes: All models are evaluated on the same 300 profiles. Machine-learning predictions are
out of fold; Qwen metrics use each profile's five-draw empirical distribution. The paired
Qwen-minus-HGB absolute-error difference is .456, with a 2,000-replication profile-cluster bootstrap
interval of $[.407,.507]$.
\end{tablenotes}
\end{threeparttable}
\end{table}

\begin{table}[!htbp]
\centering
\caption{Temporal transfer: train on 2012 and 2018, evaluate 2021}
\label{tab:s8-temporal}
\begin{threeparttable}
\small
\begin{tabular}{lrrr}
\toprule
Model & Ordinal MAE & Total variation & Variance ratio \\
\midrule
Weighted historical prior & .997 & .113 & .960 \\
Multinomial logit & .961 & .139 & .928 \\
Histogram gradient boosting & .945 & .107 & .959 \\
Qwen3-8B on matched 2021 profiles & 1.442 & .561 & .459 \\
\bottomrule
\end{tabular}
\begin{tablenotes}[flushleft]
\footnotesize
\item Notes: Every row evaluates the exact 100-profile 2021 cohort, allowing item-specific missing
human answers. The first three rows use only 2012--2018 outcomes for training; Qwen is zero-shot.
The comparison concerns reconstruction on a common target cohort, not equal training budgets.
\end{tablenotes}
\end{threeparttable}
\end{table}

\begin{table}[!htbp]
\centering
\caption{Model-call quality and software environment}
\label{tab:s9-computation}
\small
\begin{tabularx}{\textwidth}{lY}
\toprule
Component & Recorded value \\
\midrule
LLM & Qwen3-8B, non-thinking mode, local Ollama 0.31.1 \\
Decoding & Temperature .7; top-$p$ .9; fixed recorded seed \\
Hardware & Apple M4 computer; 24 GB unified memory \\
Calls & 3,331 logged attempts; 3,330 successful; success rate 99.97\% \\
Latency & Median 4.92 seconds; 90th percentile 9.66 seconds \\
Reference models & scikit-learn 1.7.2; five-fold out-of-fold evaluation \\
R environment & R 4.5.2; package versions recorded in \texttt{environment/R-session-info.txt} \\
Bootstrap & 2,000 profile-cluster replications within wave \\
Monte Carlo & 1,000 replications; $n=300$; seed 20260702 \\
Human sampling calibration & 1,000 samples per wave; $n=100$; seed 20260706 \\
\bottomrule
\end{tabularx}
\end{table}

\section{Sampling Calibration and Joint Reconstruction}

\begin{table}[!htbp]
\centering
\caption{A425 association under alternative human correlation estimators}
\label{tab:s10-correlation-sensitivity}
\begin{threeparttable}
\small
\begin{tabular}{lrrr}
\toprule
Estimator & 2012 & 2018 & 2021 \\
\midrule
Unweighted Pearson & .051 & .067 & .061 \\
Survey-weighted Pearson & .049 & .078 & .078 \\
Unweighted polychoric & .069 & .091 & .084 \\
\bottomrule
\end{tabular}
\begin{tablenotes}[flushleft]
\footnotesize
\item Notes: Entries are the mean absolute correlations between A425 and A421--A424. The weak
human association survives weighting and an ordinal latent-response estimator.
\end{tablenotes}
\end{threeparttable}
\end{table}

\begin{table}[!htbp]
\centering
\caption{Joint outcome-informed benchmark versus Qwen3-8B}
\label{tab:s11-joint-donor}
\begin{threeparttable}
\small
\begin{tabular}{lrrrrr}
\toprule
Generator & OMAE & TV & Variance ratio & Correlation RMSE & A425 mean $|r|$ \\
\midrule
Stratified joint donor & .987 & .070 & .973 & .085 & .105 \\
Qwen3-8B & 1.379 & .538 & .466 & .293 & .468 \\
\bottomrule
\end{tabular}
\begin{tablenotes}[flushleft]
\footnotesize
\item Notes: Donors are matched on wave, sex, education group, and urban residence; the matched
respondent is excluded. Five complete response vectors are sampled per profile using survey weights.
OMAE, TV, and variance use the same matched profiles as the primary benchmark. Relational metrics
compare each generated 100-profile matrix with the weighted full-sample correlation matrix.
\end{tablenotes}
\end{threeparttable}
\end{table}

\begin{table}[!htbp]
\centering
\caption{Wave-specific human sampling envelopes and Qwen3-8B estimates}
\label{tab:s12-sampling-envelope}
\begin{threeparttable}
\scriptsize
\begin{tabular}{llrrr}
\toprule
Wave & Metric & Human 2.5th & Human 97.5th & Qwen \\
\midrule
2012 & Absolute mean error & .036 & .160 & .926 \\
     & Total variation & .048 & .102 & .542 \\
     & Variance ratio & .855 & 1.137 & .442 \\
     & Correlation RMSE & .051 & .137 & .282 \\
     & A425 mean $|r|$ & .037 & .198 & .460 \\
\addlinespace
2018 & Absolute mean error & .036 & .163 & .880 \\
     & Total variation & .048 & .100 & .514 \\
     & Variance ratio & .856 & 1.141 & .592 \\
     & Correlation RMSE & .052 & .141 & .298 \\
     & A425 mean $|r|$ & .041 & .234 & .507 \\
\addlinespace
2021 & Absolute mean error & .036 & .168 & 1.021 \\
     & Total variation & .046 & .104 & .571 \\
     & Variance ratio & .846 & 1.138 & .441 \\
     & Correlation RMSE & .047 & .137 & .279 \\
     & A425 mean $|r|$ & .034 & .219 & .437 \\
\bottomrule
\end{tabular}
\begin{tablenotes}[flushleft]
\footnotesize
\item Notes: Human limits come from 1,000 repetitions of the same stratified,
probability-proportional-to-weight profile design with $n=100$ per wave. Qwen lies outside the
central 95 percent human reference envelope in all 15 cells.
\end{tablenotes}
\end{threeparttable}
\end{table}

\section{Interpretation Boundaries}
\enlargethispage{5\baselineskip}

The revised neutral condition contains five independently seeded draws per profile. The share
identical across all five draws ranges from .577 for A424 to .943 for A421. Repetition therefore
reveals nonzero stochastic variation, but the resulting total variance remains well below the human
benchmark in most wave-item cells.

The \texttt{original\_verbal} and \texttt{paraphrased} conditions contain only ten profiles each.
They are parsing diagnostics and are not used as population evidence. Only the 300-profile
\texttt{original} and \texttt{neutral\_verbal} conditions support the prompt-ablation results.

Finally, the supervised references and joint donor observe survey outcomes during training, whereas
Qwen3-8B does not. Their role is to test whether sparse profile information is a sufficient
explanation for the recorded distortion after outcome calibration. They do not isolate architecture
from calibration, establish that a supervised learner is a substitute for human respondents, or
model individual psychology.

\end{document}
